\section{References}\label{sec:references}

Sources are grouped by the role they play. Only the editions this book quotes or relies on are
listed, and each entry says what was used. Scripture is cited by book, chapter and verse in the
numbering of the Latin Vulgate. A Father's words quoted in the body with a short credit, such as
``St Leo the Great, Sermon 19'', are found here under the Father's name; where no translator is
named for a passage, the English is the book's own (Appendix~\ref{sec:scope}).

\subsection{Scripture}

\begin{itemize}[leftmargin=1.3em, itemsep=0.3em]
\item \textit{The Holy Bible}, Douay-Rheims version as revised by Richard Challoner, in the
electronic transcription of Project Gutenberg (eBook 1581), which does not name the printing it
follows. All Scripture quoted in English, unless a liturgical book or a Father is named as quoting
it.
\item \latin{Biblia Sacra Vulgatae editionis} (the Clementine Vulgate), in the electronic
transcription of eBible.org (latVUC). Scripture quoted in Latin, and the Latin behind the Douay
where the book compares the two.
\end{itemize}

\subsection{Liturgical books}

\begin{itemize}[leftmargin=1.3em, itemsep=0.3em]
\item \latin{Missale Romanum}, editio typica (Vatican City: Typis Polyglottis Vaticanis, 1962),
read in a facsimile. The stations, readings and opening words of the Masses of Advent, of the
Vigil and of Christmas, and the rubric of the Ember Saturday, pp.~1--20.
\item The Latin of the Roman Missal of 1962 as Triptych's edited text gives it (3 September and
20 October 2026): the Collects of the four Sundays of Advent, of the Ember Friday, of the Vigil of
the Nativity and of the Immaculate Conception. The edited text gives the wording of the Missal
printed at Ratisbon in 1862 (next entry), read against the typical edition of 1962.
\item \latin{Missale Romanum} (Ratisbon: Pustet, 1862). The stations of the First Sunday of
Advent and of the three Masses of Christmas, pp.~1, 16, 17 and 19; the Introit of the Vigil of the
Nativity, p.~14.
\item F.~X. Lasance and F.~A. Walsh, \textit{The New Roman Missal in Latin and English}, revised
edition (New York: Benziger, 1945). The English of the Collects: the Sundays of Advent; the Ember
Friday, p.~118; the Vigil of the Nativity, pp.~133--135; the Immaculate Conception, p.~851.
\item \latin{Breviarium Romanum}, pars hiemalis (Ratisbon: Pustet, 1888). The Latin of the Office:
the antiphon \latin{Alma Redemptoris Mater} with its versicle and prayer, pp.~130--131; the hymn
\latin{Creator alme siderum}, the chapter and the versicle \latin{Rorate} of the first Vespers of
Advent, p.~132; the first lesson of Matins and the hymns \latin{Verbum supernum prodiens} and
\latin{En clara vox redarguit}, pp.~133--137; the headings of the Isaiah lessons, pp.~133--173;
Leo's notice of the Ember days, pp.~153--154; the O antiphons and their rubric, pp.~156--157;
Matins of the Vigil of the Nativity, p.~174; the hymns \latin{Jesu Redemptor omnium} and
\latin{A solis ortus cardine}, pp.~177--178 and 189; the lessons and the responsory \latin{O magnum
mysterium} of Christmas Matins, pp.~183--187; the antiphons \latin{O admirabile commercium} and
\latin{Quando natus es} of the Circumcision, p.~231.
\item \textit{The Roman Breviary}, tr. John, Marquess of Bute, new edition, vol.~I, Winter
(Edinburgh and London: Blackwood, 1908). The English of the lessons, antiphons and responsories:
the \latin{Alma Redemptoris Mater}, pp.~210--211; \latin{Creator alme siderum} in Cardinal Newman's
English, p.~213; the Matins lessons of the Sundays and Ember Days of Advent, from Isaiah, Leo the
Great, Jerome, Gregory the Great and Ambrose, pp.~214--259; the O antiphons, pp.~243--244; the
Christmas proclamation of the Martyrology, p.~269; the lessons and fourth responsory of Christmas
Matins, pp.~274--276; the antiphons of New Year's Day, p.~315. Bute prints the Office as it stood
before 1960.
\item \latin{Martyrologium Romanum} (Rome: Propaganda Fide, 1902). The proclamation of the
Nativity, 25 December, p.~193.
\item \textit{The Little Office of the Blessed Virgin Mary}, Latin and English (New York:
Benziger, 1915). The rubric on its three forms through the year, p.~11.
\item \textit{The Gelasian Sacramentary: Liber sacramentorum Romanae ecclesiae}, ed. H.~A. Wilson
(Oxford, 1894). The five Advent Masses, pp.~214--218.
\item \textit{The Gregorian Sacramentary under Charles the Great}, ed. H.~A. Wilson (London:
Henry Bradshaw Society, 1915). The Sundays of Advent, pp.~113--116.
\item \latin{Manuale Ambrosianum}, pars altera, ed. M. Magistretti (Milan, 1904). The six Sundays
of Advent, pp.~1--41.
\item \textit{Le Liber mozarabicus sacramentorum}, ed. M. Férotin (Paris, 1912). The five Masses of
Advent, cols.~9--22.
\end{itemize}

\subsection{The prayers and hymns of Appendix~\ref{sec:prayers}}

\begin{itemize}[leftmargin=1.3em, itemsep=0.3em]
\item The Angelus: Latin and English from the Raccolta of 1910, no.~182, pp.~151--152.
\item \latin{Verbum supernum prodiens}: Latin, Breviary of 1888, p.~133; English, Caswall,
pp.~45--46.
\item \latin{En clara vox redarguit}: Latin, Breviary of 1888, pp.~136--137; English, Caswall,
pp.~46--47.
\item \latin{Creator alme siderum}: Latin, Breviary of 1888, p.~132; English, Newman's, as Bute
prints it, p.~213.
\item \latin{Alma Redemptoris Mater}: Latin, Breviary of 1888, pp.~130--131; English, Bute,
pp.~210--211.
\item \latin{Rorate caeli}: Latin and English from Guéranger, vol.~I, pp.~155--156.
\item The St Andrew prayer: Lings, p.~29.
\item The great O antiphons: Latin, Breviary of 1888, pp.~156--157; English, Bute, pp.~243--244.
\item An act of contrition: the \textit{Manual of Prayers} of 1889, p.~54.
\item \latin{Jesu Redemptor omnium}: Latin, Breviary of 1888, pp.~177--178; English, Caswall,
pp.~48--49.
\item \latin{O magnum mysterium}: Latin, Breviary of 1888, p.~183; English, Bute, p.~274.
\item \latin{A solis ortus cardine}: Latin, Breviary of 1888, p.~189; English, Caswall,
pp.~49--51.
\item \latin{O admirabile commercium}: Latin, Breviary of 1888, p.~231; English, Bute, p.~315.
\end{itemize}

\subsection{Fathers and Doctors}

\begin{itemize}[leftmargin=1.3em, itemsep=0.3em]
\item Ambrose, \latin{Expositio evangelii secundum Lucam}, ed. C. Schenkl, Corpus Scriptorum
Ecclesiasticorum Latinorum 32/4 (Vienna, 1902). Book II, 26, 41 and 42, pp.~55--56 and 63--65. His
English is Bute's, the \textit{Catena Aurea}'s, or the book's own.
\item Athanasius, \textit{On the Incarnation of the Word}, tr. A. Robertson, Nicene and Post-Nicene
Fathers, 2nd series, vol.~4 (New York, 1892). Chapter 54.3.
\item Augustine, \textit{Confessions}, tr. J.~G. Pilkington, and \textit{Letters}, tr. J.~G.
Cunningham, Nicene and Post-Nicene Fathers, 1st series, vol.~1 (Buffalo, 1886). Confessions
VIII.12.29; Letter CXXX.18.
\item Augustine, \textit{Tractates on the Gospel of John}, tr. John Gibb, Nicene and Post-Nicene
Fathers, 1st series, vol.~7 (New York, 1888). Tractate IV.1 and IV.7.
\item Augustine, \textit{Expositions on the Book of Psalms}, Nicene and Post-Nicene Fathers, 1st
series, vol.~8, in the electronic transcription of the Christian Classics Ethereal Library. The
exposition of Psalm 37 (38), paragraph 13, paraphrased.
\item Augustine, Sermons 185 and 191, in \latin{Patrologia Latina}, vol.~38 (Paris: Migne, 1845),
cols.~997--998 and 1009--1010, with the Maurist editors' note on Sermon 185.
\item Bede, \latin{In Lucae evangelium expositio}, Book I, in \latin{Patrologia Latina}, vol.~92
(Paris: Migne, 1862), col.~331. His English is the \textit{Catena Aurea}'s.
\item Bernard of Clairvaux, \textit{St Bernard's Sermons for the Seasons}, vol.~I, tr. a priest of
Mount Melleray (Dublin: Browne and Nolan, [1921]). The sermons \latin{In Adventu Domini} I, III,
IV and V, and the fourth sermon on the Virgin Mother (\latin{Super Missus est} IV).
\item Bernard of Clairvaux, \textit{Sermons of St Bernard on Advent and Christmas}, compiled and
translated at St Mary's Convent, York (London: Washbourne, 1909). The second and fourth sermons on
the Virgin Mother (\latin{Super Missus est} II.7 and IV.8) and the first sermon on the Vigil of the
Nativity, 1, 4 and 6.
\item Bernard of Clairvaux, the same sermons in \latin{Patrologia Latina}, vol.~183 (Paris: Migne,
1854), cols.~47--51 and 83--84, for the Latin quoted.
\item Cyril of Alexandria, \textit{A Commentary upon the Gospel according to S. Luke}, tr.
R. Payne Smith (Oxford: University Press, 1859), Sermon 1. The passage on Luke 2:7 that the
\textit{Catena Aurea} gives under the name ``Cyril''; it is quoted in the Catena's English.
\item Cyril of Jerusalem, \textit{Catechetical Lectures}, tr. E.~H. Gifford, Nicene and Post-Nicene
Fathers, 2nd series, vol.~7 (New York, 1894). Lecture XV, paragraphs 1, 2, 4 and 26.
\item Gregory Nazianzen, \textit{Orations}, tr. C.~G. Browne and J.~E. Swallow, Nicene and
Post-Nicene Fathers, 2nd series, vol.~7 (New York, 1894). Oration 38, on the Birthday of Christ,
1, 13 and 17; Oration 29, the Third Theological Oration, 19.
\item Gregory the Great, \latin{Homiliae in Evangelia}, in \latin{Patrologia Latina}, vol.~76
(Paris: Garnier, 1878). Homily 1.1 and 1.5, cols.~1077 and 1080; Homily 8.1, cols.~1103--1104. His
English is Bute's or the \textit{Catena Aurea}'s, except as Appendix~\ref{sec:scope} lists.
\item Jerome, \textit{Commentary on Isaiah}, Book IV, in Bute's English (Second Sunday of Advent,
fifth lesson, p.~228).
\item Leo the Great, \textit{Sermons}, tr. C.~L. Feltoe, Nicene and Post-Nicene Fathers, 2nd
series, vol.~12 (New York, 1895). Sermons XII.4, XVI.2, XVII.2, XIX.1--3, XXI.3 and LIV.4. Sermons
13 and 21 are also quoted in Bute's English.
\item Leo the Great, Sermon 13, in \latin{Patrologia Latina}, vol.~54 (Paris: Garnier, 1881),
col.~172, for \latin{Semper enim virtuti cibus jejunium fuit}.
\item Thomas Aquinas, \textit{Catena Aurea}, Oxford translation: St Luke, Part I (Oxford: Parker,
1843), on Luke 1:38, 1:46--47, 2:1--9, 3:4 and 3:5, pp.~36--112, and Part II (2nd ed., 1865), on
Luke 21:28, pp.~688--689; St John (Oxford: Parker, 1874), on John 1:23, p.~49; St Matthew (1842),
in the electronic transcription of the Christian Classics Ethereal Library, on Matthew 11:7. The
Catena is the witness for the passages given ``as the Catena gives it''.
\item Alphonsus de Liguori, \textit{The Incarnation, Birth and Infancy of Jesus Christ}, ed. Eugene
Grimm, 2nd ed. (New York: Benziger, [1886]). The contents, the editor's notice and the novena for
Christmas.
\item Francis de Sales, \textit{Introduction to a Devout Life} (New York: Pustet, undated;
translator not named). Part II, chapter 19, pp.~105--106.
\end{itemize}

\subsection{Medieval writers on the offices}

\begin{itemize}[leftmargin=1.3em, itemsep=0.3em]
\item Amalarius of Metz, \latin{Liber officialis} III.40 and \latin{Liber de ordine antiphonarii}
XXX, in \latin{Patrologia Latina}, vol.~105 (Paris: Migne), cols.~1158--1159 and 1219.
\item Honorius Augustodunensis, \latin{Sacramentarium} XCI, in \latin{Patrologia Latina}, vol.~172
(Paris: Garnier, 1895), col.~797.
\item John Beleth, \latin{Summa de ecclesiasticis officiis} LXIV, in \latin{Patrologia Latina},
vol.~202 (Paris: Migne), cols.~71--72.
\item William Durandus, \latin{Rationale divinorum officiorum}, vol.~II (Lyon, 1612), Book VI,
chapter 2, fols.~255v--256v.
\item \latin{Scala claustralium}, ascribed to Guigo the Carthusian, in \latin{Patrologia Latina},
vol.~184 (Paris: Garnier, 1879), col.~475, with the editors' headnote.
\end{itemize}

\subsection{History and councils}

\begin{itemize}[leftmargin=1.3em, itemsep=0.3em]
\item Gregory of Tours, \latin{Historia Francorum}, ed. W. Arndt, Monumenta Germaniae Historica,
\latin{Scriptores rerum Merovingicarum}, vol.~I (Hanover, 1884). Book X, chapter 31,
pp.~444--445.
\item Gregory of Tours, \textit{The History of the Franks}, tr. O.~M. Dalton, vol.~II (Oxford,
1927). Book X, chapter 31, p.~472.
\item C.~J. Hefele, \textit{A History of the Councils of the Church}, vol.~II (Edinburgh: Clark,
1876), pp.~292--293, the synod of Saragossa; vol.~IV (Edinburgh: Clark, 1895), p.~391, the council
of Tours of 567, and pp.~403--404, the first council of Mâcon, canon 9.
\item F. Maassen, \latin{Concilia aevi Merovingici}, Monumenta Germaniae Historica,
\latin{Concilia}, vol.~I (Hanover, 1893). The council of Tours of 567, canon 18, p.~126.
\item \textit{The Canons and Decrees of the Sacred and Oecumenical Council of Trent}, tr. J.
Waterworth (London: Burns and Oates, undated reprint of the translation of 1848). Session XIV,
chapters 3--4, p.~95.
\item Prosper Guéranger, \textit{The Liturgical Year: Advent}, tr. Laurence Shepherd, 2nd ed.
(Dublin: Duffy, 1870), vol.~I. The history of Advent, pp.~23--28, with the translator's note on
the English fast; the \latin{Rorate caeli}, pp.~155--156.
\end{itemize}

\subsection{Law}

\begin{itemize}[leftmargin=1.3em, itemsep=0.3em]
\item \latin{Codex Iuris Canonici} (Rome: Vatican Polyglot Press, 1918). Canon 1108 \S\,2,
pp.~316--317; canon 1252 \S\S\,1, 2 and 4, pp.~358--359.
\item \latin{Codex Iuris Canonici} (1983), Latin and English, as published on the website of the
Holy See. Canons 1246--1247 and 1250--1253.
\item Sacred Congregation of Rites, \latin{Instructio de musica sacra et sacra liturgia} (1958),
\latin{Acta Apostolicae Sedis} 50 (1958), p.~654, nn.~81 and 83.
\item Sacred Congregation of the Council, decree of 3 December 1959, \latin{Acta Apostolicae
Sedis} 51 (1959), p.~918.
\item \latin{Rubricae Breviarii et Missalis Romani}, \latin{Acta Apostolicae Sedis} 52 (1960).
Rubricae generales nn.~11, 15, 20, 24--25, 30, 71 and 131, pp.~598--620; of the Breviary
nn.~162\,a, 213, 220, 221\,a and 222\,a, pp.~625--635; of the Missal n.~299, p.~648, on the
Mass of the weekdays, and n.~507, pp.~680--681.
\item Paul VI, apostolic constitution \latin{Paenitemini}, 17 February 1966, as published on the
website of the Holy See. Norms II and VI.
\item National Conference of Catholic Bishops (United States), \textit{Pastoral Statement on
Penance and Abstinence}, 18 November 1966, as published by the United States Conference of
Catholic Bishops. Paragraphs 7, 17 and 22--25 and the note on its effective date.
\item United States Conference of Catholic Bishops, complementary norm on canons 1252--1253, with
its amendment of 1984, as published by the Conference.
\end{itemize}

\subsection{Devotional books}

\begin{itemize}[leftmargin=1.3em, itemsep=0.3em]
\item \textit{A Manual of Prayers for the Use of the Catholic Laity} (New York: Catholic
Publication Society, 1889). The list of fasting days, pp.~18--19; the act of contrition, p.~54.
\item \textit{The Raccolta}, tr. Ambrose St John (London: Burns and Oates, 1910). The Angelus,
no.~182, pp.~151--152; the Christmas novena, no.~74, p.~61.
\item A.~A. Lings, \textit{Our Favorite Novenas} (New York: Benziger, 1897). The prayer and note at
p.~29.
\item Edward Caswall, \textit{Lyra Catholica} (London: Burns, 1849). The hymns of Advent and
Christmas, pp.~45--51.
\item Ignatius of Loyola, \textit{The Spiritual Exercises}, tr. E. Mullan (New York: Kenedy, 1914).
The general examen, pp.~30--31.
\end{itemize}
